\chapter{Courbe de tonalité}
\section{Choisir un canal}
Choisir la luminance ou un canal rouge, vert ou bleu avec les boutons du module.
La diagonale représente la correspondance de référence. Une modification de
canal peut apporter une dominante colorée, même si la forme paraît proche
de celle de la courbe de luminance.

\section{Manipuler les points}
Double-cliquer dans le graphe pour ajouter une poignée, la faire glisser
pour modifier la courbe et cliquer avec le bouton droit sur une poignée
intérieure pour la supprimer. Les extrémités ne se manipulent pas comme les
points intérieurs. Éviter d’accumuler des points très proches : une courbe
plus simple est plus facile à contrôler.

Utiliser les commandes de réinitialisation de la courbe courante ou de
toutes les courbes selon l’intention. Le double-clic dans ce graphe
\emph{ajoute un point} : ce n’est pas le geste de remise à zéro des sliders.
Aucun raccourci spécifique d’ajout/suppression de points au clavier n’est déclaré.

\section{Le fond monochrome}
La silhouette grise est la luminance de l’aperçu courant ; elle reste monochrome
même si une courbe RGB est sélectionnée. C’est un repère de lecture après rendu,
pas l’histogramme exact du signal linéaire entrant dans les courbes.

\paragraph{Exemple.}
Pour donner un contraste doux, commencer par deux points intérieurs éloignés,
abaisser légèrement la zone sombre et relever légèrement la zone claire.
Comparer avec \touche{Ctrl + Maj + C} et surveiller les transitions plutôt
que chercher une forme géométrique parfaite.
